% =====================================================================================================
% Chương 3 — bản viết lại (2026-10-11). Dùng như \input{noi-dung/07-chuong3} của main.tex.
% Ghi chú cho người ghép: xem 07-chuong3-NOTES.md (nhãn ngoài chương, tài liệu mới, môi trường nhanxet).
% Mỗi dòng "% EV:" trỏ tới dòng của EVIDENCE-REGISTER.md (A01…K52), EVIDENCE-REGISTER-ADDENDUM-1.md
% (N-01…N-30, L1a-01…09) hoặc mã X3-xx (dữ kiện mã nguồn đã mở, liệt kê trong NOTES).
% =====================================================================================================

% Môi trường "Nhận xét" chưa có trong main.tex; định nghĩa có điều kiện ở đây (nên chuyển lên phần mở đầu).
\ifcsname nhanxet\endcsname\else
  \theoremstyle{definition}
  \newtheorem{nhanxet}{Nhận xét}[chapter]
\fi

\chapter{Phương pháp: luật C, lời hứa theo dự báo và sổ theo từng khách}
\label{chuong:phuongphap}

Chương này trình bày cách khóa luận đo phần dời giờ trả do điều phối cho từng khách, luật C dùng phần
đó để loại phương án, lời hứa theo dự báo, sổ theo dõi từng khách, và cách các thành phần này được cài
đặt và kiểm tra trong RideBound. Thứ tự trình bày đi từ bài toán và một lượt quyết định
(mục~\ref{muc:baitoan}), qua các định nghĩa và luật (mục~\ref{muc:bamoc}--\ref{muc:doichung}), tới lời
hứa theo dự báo, sổ theo từng khách và phần cài đặt (mục~\ref{muc:dubao}--\ref{muc:caidat}).

% -----------------------------------------------------------------------------------------------------
\section{Phát biểu bài toán và tổng quan hệ thống}
\label{muc:baitoan}

\subsection{Bài toán điều phối có lời hứa}
\label{muc:baitoan-phatbieu}

Bộ điều phối của RideBound làm việc theo từng lượt quyết định (decision epoch). Bài toán ở mỗi lượt như
sau.
\begin{itemize}[leftmargin=1.2cm,itemsep=2pt,topsep=3pt]
\item \textbf{Đầu vào:} trạng thái từng xe (vị trí, kế hoạch, sức chứa, khách trên xe và khách chờ đón);
lịch sử lời hứa của các khách đã nhận; các yêu cầu mới với điểm đón, điểm trả, khoảng giờ đón và thời
gian ngồi xe tối đa; bảng giờ đi hiện tại, được coi là giữ nguyên cho phần đường còn lại.
\item \textbf{Quyết định:} nhận hay chưa nhận từng yêu cầu mới; chọn cho mỗi xe đúng một kế hoạch, tức
danh sách điểm đón, điểm trả theo thứ tự.
\item \textbf{Đầu ra:} lời hứa giờ đón, giờ trả mới hoặc thay đổi (lời hứa đầu được công bố khi khách
xác nhận đặt xe) và một dòng sổ cho mỗi lần công bố (mục~\ref{muc:so}).
\item \textbf{Ràng buộc cứng:} cửa sổ đón, thời gian ngồi xe, sức chứa, các khóa và luật lời hứa
(Bảng~\ref{bang:rangbuoc}).
\item \textbf{Đánh giá:} lượt phục vụ, lượt trễ và tới sớm so với lời hứa đầu, độ dài thêm của lời hứa,
số khách bị điều phối dời giờ trả (mục~\ref{muc:c4-thietlap}).
\end{itemize}
% EV: X3-01, X3-02 (sự kiện; trường trạng thái xe và yêu cầu); A05 (chỉ ảnh chụp giờ đi hiện tại, giữ nguyên cho tương lai); A13; A04, X3-13 (một dòng sổ mỗi lần công bố); B03; A06; A07; A08

\subsection{Một lượt quyết định}
\label{muc:luotquyetdinh}

Hình~\ref{hinh:vongquyetdinh} và Bảng~\ref{bang:luotquyetdinh} mô tả một lượt quyết định. So với một
bộ điều phối ghép chuyến thông thường, RideBound thêm ba việc: tính các mốc giờ trả của từng khách cho
mỗi phương án (candidate plan), loại phương án vi phạm luật lời hứa trước khi gọi bộ giải (solver), và
ghi sổ sau khi chọn.
% EV: A06; X3-04 (thứ tự kiểm trong bộ kiểm quyết định); X3-13

\begin{figure}[htbp]
\centering
\begin{tikzpicture}[
  node distance=6mm,
  hop/.style={draw=xamvien,fill=xamnhat,rounded corners=3pt,align=center,
              inner sep=4pt,font=\small,minimum height=12mm,text width=2.75cm},
  mui/.style={-{Latex[length=2.2mm]},thick}
]
\node[hop] (a) {Nhận lô sự kiện\\mới};
\node[hop,right=of a] (b) {Sinh phương án\\chỉ chèn thêm};
\node[hop,right=of b] (c) {Tính ba mốc\\giờ trả cho\\từng khách};
\node[hop,right=of c] (d) {Loại phương án vi phạm ràng buộc, khóa hay trần};
\node[hop,below=9mm of d] (e) {Bộ giải chọn\\một phương án\\mỗi xe};
\node[hop,left=of e] (f) {Công bố lời hứa\\(lớp dự báo)};
\node[hop,left=of f] (g) {Ghi sổ\\theo từng khách};
\draw[mui] (a)--(b); \draw[mui] (b)--(c); \draw[mui] (c)--(d);
\draw[mui] (d)--(e); \draw[mui] (e)--(f); \draw[mui] (f)--(g);
\draw[mui,dashed] (g.west) -| node[pos=0.3,below,font=\footnotesize]{lượt sau} (a.south);
\end{tikzpicture}
\caption{Một lượt quyết định của RideBound. Ba bước tính mốc giờ trả, loại phương án theo luật lời hứa
và ghi sổ là phần thêm vào so với một bộ điều phối ghép chuyến thông thường. Trong thực nghiệm, lớp dự
báo được tính lại trên các lần chạy đã ghi (mục~\ref{muc:dubao}).}
\label{hinh:vongquyetdinh}
\end{figure}

\begin{table}[htbp]
\centering
\caption{Giả mã một lượt quyết định của RideBound}
\label{bang:luotquyetdinh}
\small
\begin{tabular}{@{}r@{\hspace{6pt}}p{9.6cm}@{\hspace{8pt}}p{4.6cm}@{}}
\toprule
\multicolumn{3}{@{}p{14.9cm}@{}}{\textbf{Đầu vào:} trạng thái đội xe; các khách đã nhận cùng lịch sử
lời hứa; lô sự kiện mới; luật lời hứa đang dùng}\\
\multicolumn{3}{@{}p{14.9cm}@{}}{\textbf{Đầu ra:} một kế hoạch cho mỗi xe; các yêu cầu mới được nhận;
lời hứa mới; các dòng sổ}\\
\midrule
1: & \textbf{hàm LượtQuyếtĐịnh}(trạng thái, lô sự kiện) & \\
2: & \quad áp lô sự kiện vào trạng thái & vị trí xe, đón, trả, giờ đi\\
3: & \quad \textbf{với mỗi} xe $k$ \textbf{thực hiện} & \\
4: & \qquad $P_k \gets \{$phương án giữ nguyên kế hoạch$\}$ & luôn dựng trước\\
5: & \qquad thêm vào $P_k$ các phương án chèn tối đa 2 yêu cầu mới, giữ thứ tự các điểm cũ; giữ
tối đa 100 phương án & bộ sinh chỉ chèn thêm\\
6: & \qquad \textbf{với mỗi} phương án $\pi \in P_k$ \textbf{thực hiện} & \\
7: & \qquad\quad \textbf{nếu} $\pi$ vi phạm cửa sổ đón, thời gian ngồi xe hay sức chứa \textbf{thì}
loại $\pi$ & ràng buộc dịch vụ\\
8: & \qquad\quad \textbf{với mỗi} khách đã nhận trên xe $k$: tính ba mốc giờ trả; \textbf{nếu} $\pi$
vi phạm khóa hay luật lời hứa \textbf{thì} loại $\pi$ & mục~\ref{muc:bamoc},
Bảng~\ref{bang:kiemtran}\\
9: & \qquad \textbf{nếu} luật lời hứa loại cả phương án giữ nguyên kế hoạch \textbf{thì} & xe bị ép;
chỉ khi mọi lỗi là lỗi hạn chót hay lỗi trần\\
10: & \qquad\quad giữ lại phương án đó làm phương án bị ép và ghi một vi phạm & luật phục hồi giữ
nguyên kế hoạch\\
11: & \qquad\quad giữ thêm các phương án bị loại không làm vi phạm đó tệ hơn & luật phục hồi không tệ
hơn\\
12: & \quad chọn cho mỗi xe đúng một phương án bằng quy hoạch nguyên & thứ tự tiêu chí ở
Bảng~\ref{bang:rangbuoc}\\
13: & \quad nhận các yêu cầu mới có trong phương án được chọn & \\
14: & \quad \textbf{với mỗi} khách có lời hứa mới hay thay đổi: công bố lời hứa, ghi một dòng sổ,
cập nhật phần đã dùng của trần & mục~\ref{muc:so}\\
15: & \quad \textbf{trả về} các kế hoạch và lời hứa & \\
16: & \textbf{kết thúc hàm} & \\
\bottomrule
\end{tabular}
\end{table}
% EV (dòng 2–5): X3-01; X3-03 (phương án giữ nguyên dựng trước; chèn giữ thứ tự điểm cũ); A24 (100 phương án, 2 yêu cầu mới mỗi xe)
% EV (dòng 6–8): X3-04; A06; A07; A10
% EV (dòng 9–11): A17; A18; X3-09 (HVA:79-99, 257-300)
% EV (dòng 12–14): A14; X3-07; A04; A06; X3-13

Luật lời hứa là một bộ lọc cứng: phương án vi phạm bị bỏ trước khi bộ giải chạy, không bị phạt trong hàm
mục tiêu. Bộ sinh phương án chỉ chèn thêm (insertion-only candidate generator) đặt điểm đón, điểm trả của
yêu cầu mới vào các cặp vị trí trong kế hoạch của xe và giữ thứ tự các điểm cũ, nên không chuyển khách đã
nhận sang xe khác. Nếu luật lời hứa loại cả phương án giữ nguyên kế hoạch (keep-route option), xe bị ép
(forced decision), và luật phục hồi (recovery rule) quyết định xe còn giữ những phương án nào
(dòng 9--11; mục~\ref{muc:doichung}).
% EV: A06; K43 (lọc cứng trước bộ giải); X3-03 (InsertionCandidateGenerator.cs:161-229, 548-575); A07 (trần 0 cho đổi xe); A17; A18

% -----------------------------------------------------------------------------------------------------
\section{Ba mốc giờ trả và phần dời do điều phối}
\label{muc:bamoc}

Giờ trả đã hứa có thể đổi vì bộ điều phối ghép thêm khách, hoặc ngay cả khi xe giữ kế hoạch cũ, vì đường
đông hơn lúc hứa hay xe đi nhanh, chậm hơn dự kiến. Để giới hạn riêng loại thứ nhất, với mỗi khách đã có
lời hứa và mỗi phương án đang xét, hệ thống so ba mốc giờ trả:
\begin{itemize}[leftmargin=1.2cm,itemsep=2pt,topsep=3pt]
\item $v$: giờ trả đã công bố cho khách ở lần gần nhất;
\item $e$: giờ trả tính lại nếu xe giữ nguyên kế hoạch cũ, xuất phát từ vị trí hiện tại với bảng giờ
đi hiện tại;
\item $p$: giờ trả nếu chọn phương án đang xét, với cùng bảng giờ đi.
\end{itemize}
Ba hiệu sau được tính bằng giây:
\begin{equation}
x = e - v, \qquad c = p - e, \qquad z = p - v = x + c .
\label{eq:bamoc}
\end{equation}
Trong đó, $x$ là phần thay đổi khi giữ nguyên kế hoạch, $c$ là phần dời do điều phối, và $z$ là phần
khách thấy; ở dòng sổ của lời hứa đầu, $v = e = p$. Luật C chỉ tính $c$. Mốc $e$ không phải ý mới: Hörl
và Zwick~\cite{horl2022} cũng tính lại giờ trả trên tuyến đang gán theo giờ đi hiện tại.
% EV: A01; A02; X3-13 (SPEC:41-50: v = e = p ở dòng đầu); J01

Mốc $e$ là giờ trả mà mô hình tính cho kế hoạch cũ trên trạng thái hiện tại, không phải giờ đến thật, nên
$x$ gồm mọi thứ đã đổi khi giữ nguyên kế hoạch: bảng giờ đi mới, tiến độ thực của xe, các sự kiện vừa xảy
ra. Trong các thực nghiệm, thời gian dừng đón, trả bằng 0 giây và tắc đường là một hệ số chung cho cả
thành phố theo khung 15 phút. Phép tách $z = x + c$ là phép tính theo mô hình, không phải phép đo nhân
quả.
% EV: A01 (lưu ý: không phải giờ đến thật, không nhận diện nhân quả); A02; X3-04 (VALID:355-361: kế hoạch cũ trên trạng thái sau lô sự kiện); B05; B06; K04

Bảng~\ref{bang:vidu-bamoc} và Hình~\ref{hinh:bamoc} là một ví dụ dựng tay: ghép khách B làm giờ trả của
khách A lùi 70 giây, nhưng chỉ 30 giây trong đó là phần dời do điều phối; 40 giây còn lại vẫn xảy ra nếu
xe không ghép ai.
% EV: ví dụ dựng tay, không phải số đo; định nghĩa theo A01

\begin{table}[htbp]
\centering
\caption{Ba mốc giờ trả của khách A trong ví dụ dựng tay của Hình~\ref{hinh:bamoc}}
\label{bang:vidu-bamoc}
\small
\begin{tabular}{@{}lp{9.6cm}r@{}}
\toprule
\textbf{Đại lượng} & \textbf{Ý nghĩa} & \textbf{Giá trị}\\
\midrule
$v$ & Giờ trả đã hứa với A ở lần công bố gần nhất & 9:10:00\\
$e$ & Giờ trả của A nếu xe giữ nguyên kế hoạch, tính lại lúc 9:00 với bảng giờ đi hiện tại & 9:10:40\\
$p$ & Giờ trả của A nếu xe ghé đón B trước & 9:11:10\\
\midrule
$x = e - v$ & Phần thay đổi khi giữ nguyên kế hoạch & $+40$ s\\
$c = p - e$ & Phần dời do điều phối & $+30$ s\\
$z = p - v$ & Phần khách thấy & $+70$ s\\
\bottomrule
\end{tabular}
\end{table}

\begin{figure}[htbp]
\centering
\begin{tikzpicture}[x=1cm,y=1cm,font=\small,
  diem/.style={draw=xamvien,fill=xamnhat,rounded corners=2pt,inner sep=3pt,minimum height=6mm},
  mui/.style={-{Latex[length=2mm]},thick}]
\node[anchor=east] at (2.6,4.0) {Giữ nguyên kế hoạch:};
\node[diem] (x1) at (3.4,4.0) {Xe};
\node[diem] (a1) at (5.6,4.0) {Trả A};
\draw[mui] (x1)--(a1);
\node[anchor=east] at (2.6,3.0) {Ghép thêm khách B:};
\node[diem] (x2) at (3.4,3.0) {Xe};
\node[diem] (b2) at (5.6,3.0) {Đón B};
\node[diem] (a2) at (7.8,3.0) {Trả A};
\node[diem] (c2) at (10.0,3.0) {Trả B};
\draw[mui] (x2)--(b2); \draw[mui] (b2)--(a2); \draw[mui] (a2)--(c2);
\draw[-{Latex[length=2mm]}] (0.8,0)--(12.2,0) node[right]{giờ};
\foreach \vitri/\tenmoc/\giomoc in {2.0/v/9:10:00, 6.8/e/9:10:40, 10.4/p/9:11:10}
  \draw (\vitri,0.12)--(\vitri,-0.12) node[below,align=center]{$\tenmoc$\\\giomoc};
\draw[mui,dashed] (2.0,0.55)--(6.8,0.55) node[midway,above]{$x = e - v = 40$ s};
\draw[mui] (6.8,0.55)--(10.4,0.55) node[midway,above]{$c = p - e = 30$ s};
\draw[mui,densely dotted] (2.0,1.45)--(10.4,1.45) node[midway,above]{$z = p - v = 70$ s};
\end{tikzpicture}
\caption{Ví dụ dựng tay: khách A được hứa trả lúc 9:10:00; lúc 9:00, xe đang chở A ghép thêm khách B.
Mũi tên nét đứt là phần thay đổi khi giữ nguyên kế hoạch ($x$), mũi tên nét liền là phần dời do điều
phối ($c$), mũi tên chấm ở trên là phần khách thấy ($z = x + c$). Luật C chỉ tính $|c|$.}
\label{hinh:bamoc}
\end{figure}

Phần dời do điều phối được tính theo độ lớn: kéo giờ trả sớm hơn $e$ hay đẩy muộn hơn $e$ 20 giây đều
tính 20 giây. Bộ sinh chỉ chèn thêm không bảo đảm $c \ge 0$ theo cấu tạo, nên định nghĩa giữ trị tuyệt
đối. Trên các lần chạy đã có, vế này không có tác dụng đo được: ở tập kiểm tra chính (34 kịch bản,
mục~\ref{muc:c4-thietlap}), mọi lần điều phối kéo giờ trả sớm hơn $e$ chỉ sớm đúng 1 mili giây.
% EV: A04; X3-05 (DELTA:139, 259-262: trị tuyệt đối); N-05 (CORRECTED: mỗi lần dời sớm đúng 1 ms; U 3 lần, C-30 4, C-120 3); ADDENDUM-1 mục 0.1–0.2 (không nêu nguyên nhân của −1 ms)

% -----------------------------------------------------------------------------------------------------
\section{Luật C: trần phần dời giờ trả do điều phối cho từng khách}
\label{muc:luatc}

\begin{dinhnghia}[Luật C]
\label{dn:luatc}
Gọi $B$ là trần (cap), tính bằng giây, $p_0$ là lời hứa đầu của khách $r$ (lời hứa không dự báo công
bố khi khách xác nhận đặt xe), và $L_r$ là tổng $|c|$ đã tính cho $r$ kể từ $p_0$, với $L_r = 0$ lúc
xác nhận. Một phương án của một xe chỉ được giữ lại nếu với mọi khách $r$ đã nhận trên xe đó, ở giai
đoạn đã nhận, chờ đón hay đang trên xe,
\begin{equation}
L_r + |p_r - e_r| \le B .
\label{eq:luatc}
\end{equation}
Khi phương án được chọn và lời hứa mới được công bố, $L_r$ được cập nhật thành $L_r + |p_r - e_r|$.
\end{dinhnghia}

Như vậy, đến hết chuyến, tổng $|c|$ của mỗi khách không vượt $B$. Có thể hình dung $B$ như một khoản ngân
sách thời gian của từng khách; $B - L_r$ là phần trần còn lại, và trần không được cấp lại sau mỗi lần dời.
Mọi khách đã nhận có cùng một mức trần, mức này không nới theo tắc đường, và phần dời được cộng qua cả
chuyến: một luật cho phép mỗi lần dời tới 30 giây sẽ cho phép ba lần dời, tổng 90 giây, còn luật C với
$B = 30$ giây thì không. Ở Bảng~\ref{bang:vidu-bamoc}, với $L_A = 0$ và $B = 30$ giây, phương án ghép B
có $|c| = 30$ giây, đúng bằng trần, nên vẫn được giữ; sau đó phần trần còn lại của A bằng 0.
% EV: A06 (BUDGET:41-105, 88: loại khi tổng > trần; LEDGER:285-288); A13; K43; E08; J05; X3-05 (BUDGET:107-114 ba giai đoạn; CFG C-30.json:17-23)

Bảng~\ref{bang:rangbuoc} tổng hợp các ràng buộc cứng và thứ tự tiêu chí của bộ chọn
(mục~\ref{muc:doichung}); Bảng~\ref{bang:kiemtran} là giả mã của phép kiểm~\eqref{eq:luatc}.

\begin{table}[htbp]
\centering
\caption{Ràng buộc cứng của một phương án và thứ tự tiêu chí của bộ chọn. Giá trị là cấu hình của các
tập thử, không phải chuẩn ngành.}
\label{bang:rangbuoc}
\small
\begin{tabular}{@{}p{2.7cm}p{7.2cm}p{4.9cm}@{}}
\toprule
\textbf{Nhóm} & \textbf{Nội dung} & \textbf{Vai trò}\\
\midrule
Ràng buộc dịch vụ & Khách được đón trong vòng 10 phút kể từ lúc gọi; thời gian ngồi xe không quá 1,5
lần thời gian đi thẳng; mỗi xe chở tối đa 4 khách & Loại phương án không khả thi trước khi xét lời
hứa\\
% EV: B03; X3-04 (VALID:201-206 kiểm vật lý trước)
Khóa sau xác nhận & Không đổi xe, điểm đón, điểm trả; giờ đón của khách đã xác nhận phải trùng giờ đón
khi giữ nguyên kế hoạch & Giữ các phần của lời hứa ngoài giờ trả; khóa so với $e$ nên tắc đường không
làm khóa bị vi phạm\\
% EV: A07; A08 (CFG C-30.json:94-98); A10; X3-06 (LOCK:42-60, 204-207)
Luật lời hứa & Luật C: $L_r + |p_r - e_r| \le B$ cho mọi khách đã nhận trên xe; hoặc một luật đối
chứng ở Bảng~\ref{bang:luat} & Bộ lọc cứng trước bộ giải\\
% EV: A06; A09; A11; K43
\midrule
Tiêu chí 1 & Ít phương án bị ép nhất (chỉ khi lượt đó có xe bị ép) & Chỉ chọn phương án bị ép khi xe
không còn lựa chọn khác\\
Tiêu chí 2 & Nhận nhiều yêu cầu mới nhất & Ưu tiên phục vụ\\
Tiêu chí 3 & Mức dùng trần lớn nhất nhỏ nhất, rồi tổng phần sửa lời hứa do điều phối trên từng chiều
(chỉ khi có trần cứng) & Ưu tiên ổn định\\
Tiêu chí 4 & Chi phí kế hoạch nhỏ nhất, tức thời lượng còn lại của tuyến dự kiến & Hiệu quả theo mô
hình, không phải chi phí vận hành thực\\
Tiêu chí 5 & Thứ tự mã phương án & Phá hòa, để kết quả lặp lại được\\
% EV: A14 (SEL:271-385); X3-07 (SEL:285-300, 323, 367, 381; HVA:612-626); A16; I09
\bottomrule
\end{tabular}
\end{table}

\begin{table}[htbp]
\centering
\caption{Giả mã phép kiểm luật C cho một phương án}
\label{bang:kiemtran}
\small
\begin{tabular}{@{}r@{\hspace{6pt}}p{9.6cm}@{\hspace{8pt}}p{4.6cm}@{}}
\toprule
\multicolumn{3}{@{}p{14.9cm}@{}}{\textbf{Đầu vào:} phương án $\pi$ của xe $k$; trần $B$; với mỗi khách
đã nhận $r$ trên xe $k$: giai đoạn và phần đã dùng $L_r$}\\
\multicolumn{3}{@{}p{14.9cm}@{}}{\textbf{Đầu ra:} giữ hay loại $\pi$; nếu loại, một bản ghi chặn}\\
\midrule
1: & \textbf{hàm KiểmTrần}($\pi$, $k$, $B$) & \\
2: & \quad \textbf{với mỗi} khách đã nhận $r$ trên xe $k$, theo thứ tự mã khách \textbf{thực hiện} & \\
3: & \qquad $e_r \gets$ giờ trả của $r$ khi chạy lại kế hoạch cũ của $k$ từ vị trí hiện tại, với bảng giờ
đi hiện tại & mốc giữ nguyên kế hoạch\\
4: & \qquad $p_r \gets$ giờ trả của $r$ theo $\pi$, với cùng bảng giờ đi & \\
5: & \qquad $c_r \gets p_r - e_r$ & phần dời do điều phối\\
6: & \qquad \textbf{nếu} $r$ đang ở giai đoạn đã nhận, chờ đón hay trên xe \textbf{và}
$L_r + |c_r| > B$ \textbf{thì} & sớm hay muộn đều tính\\
7: & \qquad\quad \textbf{trả về} loại $\pi$, kèm bản ghi ($r$, $L_r$, $|c_r|$, $L_r + |c_r|$, $B$) &
dừng ở khách vượt trần gặp trước\\
8: & \quad \textbf{trả về} giữ $\pi$ & nếu $\pi$ được chọn: $L_r \gets L_r + |c_r|$\\
9: & \textbf{kết thúc hàm} & \\
\bottomrule
\end{tabular}
\end{table}
% EV (dòng 2–5): X3-04 (VALID:233-239 thứ tự mã khách trên xe đang xét; HVA:171; VALID:355-361 mốc e; VALID:494-499)
% EV (dòng 6–8): X3-05 (BUDGET:72-104, 107-114; DELTA:139); A06; A23 (dừng ở lỗi gặp trước); BUDGET:6-12 (trường của bản ghi chặn)

Phép kiểm dừng ở khách vượt trần gặp trước, nên mỗi bản ghi chặn cho một lý do loại phương án, không phải
mọi lý do. Cài đặt thực kiểm cùng lúc mười chiều cam kết của mỗi khách; trong các cấu hình luật C, giờ trả
là chiều có trần dương, còn đổi xe, đổi điểm đón và đổi điểm trả có trần 0.
% EV: A23; I08; A04 (10 chiều; DIM:3-32); A07; CFG C-30.json:7-89

\subsection{Phương án giữ nguyên kế hoạch dưới luật C}
\label{muc:giunguyen}

Nếu luật lời hứa loại hết phương án của một xe, kể cả phương án giữ nguyên kế hoạch, xe bị ép và lời hứa
của ít nhất một khách bị phá. Gọi \textbf{(N)} là giả thiết: ở lượt đang xét, phương án giữ nguyên kế
hoạch được bộ sinh dựng ra và qua mọi kiểm tra không thuộc luật C (ràng buộc dịch vụ và khóa).
% EV: A17; A18 (bị ép = giữ kế hoạch cũ và ghi vi phạm); GLOSSARY-VI mục E

\begin{nhanxet}
\label{nx:giunguyen}
Giả sử (N) và mọi khách đã nhận trên xe có $0 \le L_r \le B$. Khi đó luật C không loại phương án giữ
nguyên kế hoạch.
\end{nhanxet}

\noindent\textit{Chứng minh.} Phương án giữ nguyên dùng đúng kế hoạch cũ, được tính từ cùng trạng thái
và cùng bảng giờ đi như mốc $e$, nên $p_r = e_r$ và $c = 0$ với mọi khách trên xe. Vì vậy
$L_r + |p_r - e_r| = L_r \le B$ và phép kiểm~\eqref{eq:luatc} cho qua; theo (N), phương án cũng qua mọi
kiểm tra khác. Điều kiện $L_r \le B$ được giữ qua các lượt vì $L_r$ bắt đầu từ 0 và chỉ tăng khi một
phương án đã qua~\eqref{eq:luatc} được chọn. \hfill$\square$
% EV: A22 (theo cấu tạo, giả thiết N); A06; X3-03; X3-04

Nhận xét đúng theo cấu tạo, không phải kết quả thực nghiệm; dữ liệu chỉ khớp với nó. Trên tập kiểm tra
chính, luật C có 0 quyết định bị ép trong 20.101 quyết định với trần 30 giây và 0 trong 20.435 quyết định
với trần 120 giây.
% EV: A21; N-21 (C-30 0 quyết định bị ép); X3-16 (STG:43: 20.101 quyết định); N-07 (0 / 20.435)

Nhận xét có ba giới hạn. Thứ nhất, nó không bảo đảm xe luôn có phương án hợp lệ: nếu (N) không thỏa,
chẳng hạn khi kế hoạch đang chạy đã vi phạm một ràng buộc vật lý, xe vẫn có thể kẹt, nhưng do ràng buộc
khác. Thứ hai, $c = 0$ không có nghĩa là khách không thấy thay đổi: với $v = 600$ và $e = p = 660$
(giây), phương án giữ nguyên có $z = 60$, nên luật C không bảo đảm giờ đến thật. Thứ ba, một lần sửa giữ
đúng lời hứa vẫn bị tính: với $v = p = 600$ và $e = 720$, phương án kéo giờ trả về 600 giây có $z = 0$
nhưng $|c| = 120$. Mục~\ref{muc:doichung} và Phụ lục~\ref{phuluc:B} so tính chất của
Nhận xét~\ref{nx:giunguyen} với các luật khác.
% EV: A22 (giả thiết N); X3-03 (bộ sinh trả lỗi khi kế hoạch đang chạy vi phạm ràng buộc vật lý); K05; TEX/07-chuong3.tex:650-654 (CE-2, CE-3, ví dụ dựng tay)

% -----------------------------------------------------------------------------------------------------
\section{Các luật đối chứng và bộ chọn phương án}
\label{muc:doichung}

Chương~\ref{chuong:thucnghiem} so luật C với bốn luật khác chạy trên cùng bộ điều phối.
Bảng~\ref{bang:luat} định nghĩa chúng một lần; ký hiệu ở cột thứ hai chỉ dùng trong các bảng.

\begin{table}[htbp]
\centering
\caption{Các luật được so sánh. Cột thứ hai là ký hiệu chỉ dùng trong bảng; $B$ là trần (với luật hạn
chót là độ chùng), tính bằng giây; $p_0$ là lời hứa đầu.}
\label{bang:luat}
\footnotesize
\setlength{\tabcolsep}{3pt}
\begin{tabular}{@{}p{2.1cm}p{1.1cm}p{3.4cm}p{2.2cm}p{2.3cm}p{3.6cm}@{}}
\toprule
\textbf{Tên trong bài} & \textbf{Ký hiệu} & \textbf{Kiểm điều gì} & \textbf{Khóa} & \textbf{Khi xe bị ép}
& \textbf{Dùng ở tập nào}\\
\midrule
Điều phối không trần & U & Không kiểm giờ trả & Xe, điểm đón, giờ đón; không đổi điểm trả & Không xảy
ra, vì khóa so với $e$ & Tập kiểm tra chính; tập chỉnh trần buổi sáng; tập tắc đường cố định (tên cũ
LOCKS)\\
% EV: A07; A08; K50; A10; X3-16 (STG:42: 0 quyết định bị ép); B10; B15; B17
Ghép tự do & FREE & Không kiểm; chính sách chi phí cuốn chiếu mặc định & Không khóa & Không áp dụng &
Tập tắc đường cố định\\
% EV: A17; K50; B17
Luật C & C-$B$ & $\sum |c| \le B$ cả chuyến, cho mọi khách đã nhận & Như U & Không xảy ra theo cấu tạo
(Nhận xét~\ref{nx:giunguyen}) & Tập chính ($B$ = 30, 120); buổi sáng (30--300); tắc đường cố định (30,
60, 120); chiều tuyết (30)\\
% EV: A06; A22; B10; L1a-01; D01; B15; B16; B17; B20; N-26 (chiều tuyết 0 / 5.067 quyết định bị ép)
Luật hạn chót & M$^{+}$-$B$ & $p - p_0 \le B$ ở mọi lần cập nhật; chỉ chặn phía muộn & Như U & Phục hồi
không tệ hơn & Tập chính (30, 60); buổi sáng (30--300)\\
% EV: A09; A19; B10; D04; B15; B16
Luật trần độ dời khách thấy & V-$B$ & $\sum |z| \le B$ cả chuyến & Như U & Phục hồi không tệ hơn & Tập
chính (30); buổi sáng (30)\\
% EV: A11; A19; B10; B15
\bottomrule
\end{tabular}
\end{table}

Luật hạn chót chỉ mượn phần hạn chót của Schulz và Pfeiffer~\cite{schulz2026}, ở mức 30--60 giây thay
vì 900 giây, không tái lập thuật toán của họ. Trên tập tắc đường cố định, điều phối không trần mang tên
cũ LOCKS và chạy bằng bộ chạy cũ không có luật phục hồi. Phụ lục~\ref{phuluc:A} gọi cấu hình không khóa,
không trần của tập kịch bản tĩnh là B1, và cấu hình có khóa cùng trần 30 giây là C1.
% EV: J10; J11; B17; K33; X3-18 (TEX/08-chuong4.tex:128-136)

Ba luật lời hứa khác nhau ở đại lượng được cộng. Khi mỗi lần cập nhật so với lời hứa ngay trước nó
($v_t = p_{t-1}$), sau $k$ lần cập nhật
\begin{equation}
\sum_{t=1}^{k} (x_t + c_t) \;=\; \sum_{t=1}^{k} z_t \;=\; p_k - p_0 .
\label{eq:dongnhat}
\end{equation}
Luật hạn chót kiểm $p_k - p_0 \le B$, tức cộng có dấu cả hai phần; luật trần độ dời khách thấy cộng
$|z_t|$; luật C chỉ cộng $|c_t|$. Vì vậy riêng tắc đường có thể làm hai luật đầu loại phương án giữ
nguyên: ở Bảng~\ref{bang:vidu-bamoc}, nếu 9:10:00 là lời hứa đầu, phương án giữ nguyên có
$e - p_0 = |z| = 40$ giây, vượt trần 30 giây của cả hai luật, dù $c = 0$.
% EV: A12 (SPEC:50, 62-63: chuỗi liên tục); A09; A10 (LOCK:156-161); A11

Gần nhất với cách đo này là Hörl và Zwick~\cite{horl2022}: ở mỗi lượt quyết định, họ duyệt tuyến đang
gán theo giờ đi hiện tại để có giờ trả cập nhật, tức mốc $e$. Công thức 7 của họ đặt giờ trả bắt buộc của
khách trên xe đang xét bằng giá trị lớn hơn giữa giờ cập nhật và giờ trả muộn nhất, kèm phạt cho mỗi giây
vượt giờ trả muộn nhất. Với giờ trả muộn nhất là $p_0 + B$, đó là hạn chót nới $\max(p_0 + B, e)$, một
cách viết lại của khóa luận chứ không phải định nghĩa trong bài. Hạn chót nới này cũng có tính chất của
Nhận xét~\ref{nx:giunguyen} vì phương án giữ nguyên có $p = e$; đây là suy luận, chưa được chạy. Luật C
khác ở ba chỗ: một trần chung cho mọi khách đã nhận, không nới theo tắc đường, và cộng qua cả chuyến.
% EV: J01 (HZ:529-543); J03 (HZ:562-606: công thức 7, ràng buộc 9–10, phạt 11); J05; A22 (lưu ý: suy luận V17 R6); K02; K26 (không dùng "cứng thay vì phạt" làm khác biệt); K27

Bộ chọn phương án (selector) chọn cho mỗi xe đúng một phương án bằng quy hoạch nguyên, giải bằng
OR-Tools~\cite{ortools}, theo thứ tự tiêu chí (lexicographic order) của Bảng~\ref{bang:rangbuoc}: chỉ khi
hòa ở tiêu chí trước mới xét tiêu chí sau. Tiêu chí 1 chỉ có mặt khi có xe bị ép. Tiêu chí 3 bật cả khi
chỉ có các trần 0 của khóa, nên có ở cả điều phối không trần; phần đầu của nó không phân biệt phương án
vì chiều có trần 0 luôn được tính là đã dùng hết. Trên tập kiểm tra chính, mức này bằng 1.000.000 phần
triệu ở mọi lượt có mức này: 3.503 lượt khi không trần, 3.506 lượt với trần 30 giây và 3.504 lượt với
trần 120 giây.
% EV: A14; X3-07 (SEL:276-280, 285-287); A15 (HVA:612-626; U 3.503/3.503, C-30 3.506/3.506); N-25 (C-120 3.504/3.504); A16

Điều phối không trần và luật C dùng cùng tệp chính sách và cùng bộ chọn, chỉ khác ở trần giờ trả; hiệu
giữa chúng là phép so của việc thêm trần khi giữ nguyên bộ chọn. Kết quả có còn khi bộ chọn ưu tiên chi
phí hay không thì chưa được kiểm (mục~\ref{muc:c4-vaitro}).
% EV: A08 (diff U.json / C-30.json: chỉ hardLimit 30000 của giờ trả); H01; H02; K18

Có hai luật phục hồi. Luật giữ nguyên kế hoạch giữ phương án giữ nguyên làm phương án bị ép và ghi một
vi phạm, nếu mọi lỗi của nó là lỗi hạn chót hay lỗi trần. Luật không tệ hơn kế hoạch giữ còn giữ, trên xe
đã bị ép, các phương án mắc cùng lỗi nhưng không làm lỗi đó tệ hơn, để xe vẫn nhận được khách mới; khóa
xe và khóa theo giai đoạn không bao giờ được miễn, và mỗi lần miễn được ghi thành vi phạm. Tập kiểm tra
chính và tập chỉnh trần buổi sáng dùng luật thứ hai cho mọi cấu hình.
% EV: A17; A18 (ADR:5066-5092; HVA:85-99, 257-300); X3-09; A19 (RUN38:9; WP4:7); X3-10 (run_w07.py:8; run_w07x.py:7; run_t38c120.py:3-8)

% -----------------------------------------------------------------------------------------------------
\section{Lời hứa theo dự báo}
\label{muc:dubao}

Lời hứa của bộ điều phối giả định bảng giờ đi lúc công bố giữ nguyên, nên khi đường đông dần, lời hứa
sớm hơn giờ đến thật dù không có quyết định điều phối nào. Lời hứa theo dự báo (forecast-aware promise),
một cách làm đã có trong thực tiễn (mục~\ref{muc:c2-dubao}), cộng sẵn mức tăng tắc đường mà các ngày
trước cho thấy ở cùng khung giờ. Với lúc công bố $t$ và lời hứa không dự báo $p$, lời hứa kéo dài $S$ là
nghiệm của
\begin{equation}
\int_{t}^{S} \frac{\mathrm{d}\tau}{g(\tau)} = p - t, \qquad
g(\tau) =
\begin{cases}
1 & \text{nếu } b(\tau) = b_0,\\
m(b(\tau))/m(b_0) & \text{nếu không,}
\end{cases}
\label{eq:keodan}
\end{equation}
và lời hứa hiển thị là $D = \max(p, S)$. Trong đó, $m(b)$ là trung bình hệ số tắc đường ở khung 15 phút
$b$ trên các ngày nguồn, $b(\tau)$ là khung chứa thời điểm $\tau$, và $b_0 = b(t)$: mức hiện tại lấy từ
bảng giờ đi đang dùng, còn từ các ngày khác chỉ lấy hình dạng tăng, giảm. Phép lấy $\max$ làm lời hứa
không bao giờ bị rút ngắn (giả mã ở Bảng~\ref{bang:dubao}).
% EV: A05; F01 (DISP:1-8, 60-64; STRETCH:1-12, 25-52); F10 (kéo dãn theo hệ số thật khớp trong 2 s với khách có phần điều phối bằng 0); B06; J35; J21; X3-11

Ngày nguồn là các ngày cùng loại (ngày thường hay cuối tuần) trước ngày của kịch bản, không gồm ngày có
tuyết 15/11/2018, và cần ít nhất 2 ngày; vì vậy trong dữ liệu 12--18/11/2018 chỉ các ngày 14, 15 và 16/11
có dự báo. Trên tập kiểm tra chính, 21 trong 34 kịch bản có dự báo; 13 kịch bản ngày 17 và 18/11 giữ lời
hứa không dự báo.
% EV: F02 (DISP:10-16, 34-52); F04 (21 kịch bản ngày thường có dự báo); N-16, N-28 (13 kịch bản cuối tuần không có dự báo)

\begin{table}[htbp]
\centering
\caption{Giả mã lời hứa theo dự báo cho một lời hứa đã công bố}
\label{bang:dubao}
\small
\begin{tabular}{@{}r@{\hspace{6pt}}p{9.6cm}@{\hspace{8pt}}p{4.6cm}@{}}
\toprule
\multicolumn{3}{@{}p{14.9cm}@{}}{\textbf{Đầu vào:} lúc công bố $t$; lời hứa không dự báo $p$; ngày $d$
của kịch bản; hệ số tắc đường theo khung 15 phút của các ngày trước $d$}\\
\multicolumn{3}{@{}p{14.9cm}@{}}{\textbf{Đầu ra:} lời hứa hiển thị $D$}\\
\midrule
1: & \textbf{hàm LờiHứaTheoDựBáo}($t$, $p$, $d$) & \\
2: & \quad $N \gets$ các ngày cùng loại với $d$ và trước $d$, bỏ ngày 15/11 & chỉ dùng ngày đã qua\\
3: & \quad \textbf{nếu} $N$ có ít hơn 2 ngày \textbf{hoặc} $p \le t$ \textbf{thì trả về} $p$ & giữ lời
hứa không dự báo\\
4: & \quad $m(b) \gets$ trung bình hệ số của khung $b$ trên các ngày trong $N$ & \\
5: & \quad $W \gets p - t$;\; $\tau \gets t$;\; $b_0 \gets b(t)$ & thời gian còn lại theo tốc độ hiện
tại\\
6: & \quad \textbf{lặp} & \\
7: & \qquad $g \gets 1$ nếu $b(\tau) = b_0$, ngược lại $g \gets m(b(\tau))/m(b_0)$ & mức tắc đường so
với hiện tại\\
8: & \qquad $h \gets$ thời gian từ $\tau$ tới hết khung $b(\tau)$ & \\
9: & \qquad \textbf{nếu} $h/g \ge W$ \textbf{thì} $S \gets \tau + W g$ và thoát vòng lặp & \\
10: & \qquad $W \gets W - h/g$;\; $\tau \gets$ đầu khung kế tiếp & \\
11: & \quad \textbf{trả về} $D \gets \max(p, S)$ & không rút ngắn lời hứa\\
12: & \textbf{kết thúc hàm} & \\
\bottomrule
\end{tabular}
\end{table}
% EV (dòng 2–4): F02; X3-11 (DISP:34-57)
% EV (dòng 5–11): F01; X3-11 (STRETCH:25-39, 49-52)

Trong thực nghiệm, lời hứa theo dự báo được tính lại trên các lần chạy đã ghi, gọi là phát lại (replay):
quyết định điều phối không đổi, chỉ lời hứa đổi. Ký hiệu $D_0$ là lời hứa đầu có dự báo; giá của dự báo
là độ dài thêm $D_0 - p_0$. Còn một chỗ chưa khớp: phép kiểm khả thi vẫn dùng giờ không dự báo, nên dưới
luật C với trần 30 giây, 34 trong 1.471 khách được phục vụ của 21 kịch bản có dự báo có lời hứa giờ đón
hiển thị muộn hơn giờ đón muộn nhất; chỗ này phải xử lý trước khi triển khai. Phép thử đưa dự báo vào
chính bộ điều phối được trình bày ở mục~\ref{muc:c4-dubao}.
% EV: F03 (chỉ phần "phát lại trên các lần chạy đã ghi, không chạy mô phỏng mới"); F01; F06; X3-12 (FFA:86, 93: D0 = lời hứa hiển thị đầu); F07 (FR:70); F12

% -----------------------------------------------------------------------------------------------------
\section{Sổ theo dõi từng khách}
\label{muc:so}

Sổ theo từng khách (per-rider ledger) do RideBound ghi trong lúc chạy: mỗi khách có một dòng cho lời hứa
đầu và một dòng cho mỗi lần cập nhật, gồm $v$, $e$, $p$, các độ lớn $|x|$, $|c|$, $|z|$ trên mười chiều
cam kết và phần trần đã dùng $L$; trạng thái cuối của mỗi lần chạy chứa toàn bộ sổ. Sau khi chạy, một
phép tính trên sổ chia độ trễ của mỗi khách so với lời hứa đầu thành ba phần:
\begin{equation}
a - p_0 \;=\; q \;+\; f \;+\; (a - D_0 - q), \qquad q = \sum_{t} c_t, \qquad f = D_0 - p_0 .
\label{eq:so}
\end{equation}
% EV: A04 (DELTA:64-97, 101-160; DIM:3-32); A06 (LEDGER:285-288); X3-13 (SPEC:35-46); X3-12 (FFA:91, 134-138)

Trong đó, $a$ là giờ trả thực tế; $q$ là \textit{phần tính cho điều phối}, một tổng có dấu, khác với
$\sum|c|$ mà luật C giới hạn (dưới luật C, $|q| \le \sum|c| \le B$); $f$ là \textit{phần được hồ sơ lịch
sử giải thích}, bằng 0 khi không có dự báo; số hạng cuối là \textit{phần còn lại}. Khi chuỗi cập nhật
liên tục, theo~\eqref{eq:dongnhat} phần còn lại bằng $\sum_t x_t + \rho - f$, với $\rho = a - p_k$ là độ
lệch giữa giờ trả thật và lời hứa cuối $p_k$. Chương~\ref{chuong:thucnghiem} cộng ba phần này trên các
khách trễ hơn lời hứa đầu quá 60 giây.
% EV: X3-12 (FFA:91 q = Σ(p − e) có dấu; FFA:134-138 ba phần trên khách trễ > 60 s); A12; A06; F09

Đây là phép tính theo mô hình, không kết luận ai có lỗi: tên ``phần tính cho điều phối'' chỉ nói rằng mô
hình gán phần này cho quyết định điều phối khi so với kế hoạch giữ nguyên. Phần được hồ sơ lịch sử giải
thích phụ thuộc vào dự báo đang dùng, và phần còn lại là phần dự báo đó chưa giải thích được, không phải
phần ``không thể biết trước''.
% EV: K04; A01 (lưu ý D03:222-224); F09 (lưu ý); K19

% -----------------------------------------------------------------------------------------------------
\section{Cài đặt, kiểm thử và tái lập}
\label{muc:caidat}

\subsection{Kiến trúc và trao đổi với bộ mô phỏng}
\label{muc:kientruc}

RideBound được viết bằng C\# trên .NET 10, chia thành các dự án theo tầng, phụ thuộc từ ngoài vào trong
(Bảng~\ref{bang:kientruc}). Tầng nghiệp vụ (Domain) không tham chiếu dự án hay thư viện nào, tầng ứng dụng
(Application) chỉ tham chiếu tầng nghiệp vụ, OR-Tools chỉ được tham chiếu trong thành phần cầu nối với bộ
giải, và một bộ kiểm thử kiến trúc kiểm các quy tắc này. Nhờ vậy phép kiểm trần, ba mốc giờ trả và sổ
không phụ thuộc vào bộ mô phỏng hay bộ giải.
% EV: X3-14 (Directory.Build.props net10.0; các tệp .csproj; DependencyRuleTests.cs:11-29, 95-110, 137-166)

\begin{table}[htbp]
\centering
\caption{Các thành phần chính của RideBound và bộ chuyển đổi FleetPy}
\label{bang:kientruc}
\small
\begin{tabular}{@{}p{4.0cm}p{10.8cm}@{}}
\toprule
\textbf{Thành phần} & \textbf{Vai trò}\\
\midrule
Domain & Mô hình nghiệp vụ: yêu cầu và giai đoạn của khách, lời hứa, khóa, sổ theo từng khách, phép
kiểm trần. Không tham chiếu dự án hay thư viện nào khác.\\
Application & Tính lại kế hoạch theo bảng giờ đi, tính ba mốc giờ trả và kiểm một phương án theo khóa
và luật lời hứa.\\
Algorithms & Bộ sinh phương án chỉ chèn thêm, đánh giá phương án theo luật lời hứa và luật phục hồi,
bộ chọn phương án.\\
Solvers.OrTools & Cầu nối tới bộ giải quy hoạch nguyên của OR-Tools 9.15.\\
Contracts & Định dạng các bản tin trao đổi với bộ mô phỏng.\\
Runner & Tiến trình nhận bản tin, ghép các thành phần trên và trả quyết định.\\
Benchmarking & Chuẩn hóa dữ liệu Manhattan của FleetPy thành kịch bản thử.\\
Bộ chuyển đổi FleetPy (Python) & Đóng vai trò phần điều khiển đội xe của FleetPy 1.0.2: gửi sự kiện
cho bộ chạy và áp quyết định vào mô phỏng.\\
\bottomrule
\end{tabular}
\end{table}
% EV (Domain, Application, Algorithms): X3-14; BUDGET, LEDGER, LOCK, POLICY (Domain); VALID, DELTA, PROJ, PPROJ (Application); SEL, HVA, InsertionCandidateGenerator.cs (Algorithms)
% EV (Solvers.OrTools, Contracts, Runner, Benchmarking, bộ chuyển đổi): X3-14 (Google.OrTools 9.15.6755); X3-15; NORM; B01

Bộ mô phỏng là FleetPy 1.0.2~\cite{engelhardt2026} với bộ dữ liệu chuẩn Manhattan. Một bộ chuyển đổi
viết bằng Python đóng vai trò phần điều khiển đội xe của FleetPy; nó khởi động bộ chạy RideBound như một
tiến trình con và trao đổi qua luồng vào, ra chuẩn, mỗi bản tin một dòng JSON (Hình~\ref{hinh:bantin}).
Bộ chạy từ chối lô sự kiện mới khi quyết định trước chưa được báo đã áp dụng, nên các quyết định được áp
dụng đúng thứ tự. Các bản tin của mỗi lần chạy được lưu lại, và các con số ở
Chương~\ref{chuong:thucnghiem} được tính từ bản ghi này.
% EV: B01; J18 (FLEET:474: FleetPy 1.0.2); X3-15 (README:1-12; runner_client.py:133-135; Program.cs:104-105; NdjsonReader/Writer; RunnerSession.cs:203-207, 448-455, 543); X3-13 (SPEC:35-37); A26 (T38I: đọc transcript-00.ndjson)

\begin{figure}[htbp]
\centering
\begin{tikzpicture}[y=0.75cm,font=\small,
  ben/.style={draw=xamvien,fill=xamnhat,rounded corners=2pt,align=center,inner sep=4pt,text width=4.3cm},
  tin/.style={-{Latex[length=2mm]},thick},
  nhan/.style={midway,above,font=\footnotesize,align=center}]
\node[ben] (L) at (0,0) {FleetPy 1.0.2\\và bộ chuyển đổi (Python)};
\node[ben] (R) at (9.5,0) {Bộ chạy RideBound\\(C\#, .NET 10)};
\draw[xamvien,thick] (L.south)--(0,-10.3);
\draw[xamvien,thick] (R.south)--(9.5,-10.3);
\draw[tin] (0,-1.3)--(9.5,-1.3) node[nhan]{chào hỏi; khởi tạo lần chạy (kịch bản, cấu hình luật)};
\draw[tin] (9.5,-2.0)--(0,-2.0) node[nhan]{đã khởi tạo};
\draw[dashed,xamvien] (-0.6,-2.5) rectangle (12.2,-7.2);
\node[anchor=south west,font=\footnotesize\itshape] at (-0.55,-7.15) {mỗi lượt quyết định};
\draw[tin] (0,-3.6)--(9.5,-3.6) node[nhan]{lô sự kiện: yêu cầu mới, xác nhận đặt xe,\\tiến độ xe, đón,
trả, bảng giờ đi mới};
\node[ben,text width=3.9cm,font=\footnotesize] at (9.5,-4.6) {sinh phương án, kiểm luật, chọn, ghi sổ
(Bảng~\ref{bang:luotquyetdinh})};
\draw[tin] (9.5,-5.8)--(0,-5.8) node[nhan]{quyết định: kế hoạch mới của các xe, lời hứa};
\draw[tin] (0,-6.5)--(9.5,-6.5) node[nhan]{đã áp dụng quyết định};
\draw[tin] (0,-7.9)--(9.5,-7.9) node[nhan]{yêu cầu trạng thái};
\draw[tin] (9.5,-8.6)--(0,-8.6) node[nhan]{trạng thái, gồm sổ theo từng khách};
\draw[tin] (0,-9.3)--(9.5,-9.3) node[nhan]{kết thúc lần chạy};
\draw[tin] (9.5,-10.0)--(0,-10.0) node[nhan]{tóm tắt lần chạy};
\end{tikzpicture}
\caption{Các bản tin giữa bộ mô phỏng FleetPy và bộ chạy RideBound trong một lần chạy. Mỗi bản tin là một
dòng JSON trên luồng vào, ra chuẩn; phần trong khung nét đứt lặp lại ở mỗi lượt quyết định.}
\label{hinh:bantin}
\end{figure}
% EV: X3-15 (ProtocolPrimitives.cs:104-124: hello, initializeRun, eventBatch, decision, decisionApplied, checkpoint, finalizeRun, runSummary; EventBatchMessages.cs:42-61); X3-13

\subsection{Giới hạn tính toán và thời gian chạy}
\label{muc:gioihan}

Bộ sinh có giới hạn: tối đa 100 phương án và 2 yêu cầu mới mỗi xe trong một lượt, cùng các giới hạn khối
lượng cho bước sinh, bước kiểm và bộ giải (OR-Tools 9.15, hạt giống cố định 7), nên không vét cạn;
mục~\ref{muc:c4-hople} báo các bộ đếm cho biết các giới hạn này có cắt bớt phương án hay không. Phép kiểm
ở Bảng~\ref{bang:kiemtran} duyệt các khách đã nhận trên xe và so mười chiều cam kết cho mỗi phương án, nên
khối lượng kiểm ở một xe tỉ lệ với số phương án nhân số khách trên xe, nhân độ dài kế hoạch.
% EV: A24 (WP4:9-23); X3-17; N-25; X3-04 (VALID:233-239 vòng khách trên xe; 355-361 tính lại kế hoạch cũ mỗi khách; 494-520); X3-05 (BUDGET:76-99 vòng 10 chiều)

Thời gian chạy chỉ được đo cho cả lần chạy: trên tập tắc đường cố định (16 lần chạy mỗi cấu hình, bộ chạy
cũ), trung vị là 730 giây với ghép tự do, 667 giây với điều phối không trần và 641 giây với luật C trần
30 giây. Chi phí thêm của phép kiểm trần ở từng lượt quyết định chưa được đo riêng.
% EV: H13 (EL17:2-6); X3-19 (b1-ref = FREE, c1-nobudget = LOCKS, c1-h6ref = C-30); B17

\subsection{Kiểm tra kết quả}
\label{muc:kiemtra}

Kết quả được kiểm ở bốn mức.
\begin{itemize}[leftmargin=1.2cm,itemsep=2pt,topsep=3pt]
\item \textbf{Kiểm thử phần mềm:} khi luật phục hồi không tệ hơn được thêm vào mã, bộ kiểm thử của
RideBound qua 1.025 trên 1.025 trường hợp; đây là số ghi lại lúc đó, không được chạy lại khi viết khóa
luận.
\item \textbf{Kiểm đặc tả:} ở cả 382.006 lần cập nhật giờ trả của 185 lần chạy (170 lần chạy của năm luật
trên tập kiểm tra chính và 15 lần chạy hiệu chỉnh cùng đợt), $|x|$, $|c|$, $|z|$ tính lại từ ba mốc khớp
đúng các trường ghi trong sổ.
\item \textbf{Tính lại độc lập:} một chương trình viết riêng từ đặc tả chỉ số tính lại 13 hạng mục kết quả
từ bản ghi thô, không hạng mục nào lệch; các số đếm chính trên tập kiểm tra chính được ba chương trình đọc
bản ghi tính lại và khớp nhau. Các chương trình kiểm ở mức này và mức trên do các mô hình ngôn ngữ lớn
Claude và Codex viết và chạy, không phải kiểm thủ công; chúng cho biết con số khớp với bản ghi, không cho
biết mô phỏng giống thực tế đến đâu.
\item \textbf{Tái lập:} tập kiểm tra chính và tập chỉnh trần buổi sáng dùng cùng một bản chạy đã lưu trữ,
có luật phục hồi; tập tắc đường cố định dùng một bản cũ hơn, không có luật phục hồi. Một bản dựng lại từ
mã nguồn phải cho đúng kết quả cũ trên một số kịch bản trước khi được dùng (Phụ lục~\ref{phuluc:C}).
\end{itemize}
% EV (kiểm thử): A27 (ADR:5100, số lịch sử)
% EV (đặc tả): A26 (T38:106; T38I:68); B10 (170 + 15)
% EV (tính lại): C19 (T38I:79-97); N-03, N-06, N-07 (3 bộ đọc khớp); C04, E01 (hai bộ đọc độc lập); A26 lưu ý ("agent"); khóa A, FJC, FJS, STG, CXS, ADD1 (Codex / Claude); K41; I01
% EV (tái lập): A28 (RUNW07:40-41); X3-10 (run_w07.py:8; run_t38c120.py:3-8); B17; I13

Chương~\ref{chuong:thucnghiem} dùng các thành phần trên để đo luật C được gì và mất gì trên mô phỏng
FleetPy, so với điều phối không trần, luật hạn chót và cách đệm lời hứa.
